% Provisional source-preserving import from Person A draft snapshot.
% Source: tmp/docs/personA_submission_extracted.txt
% Expanded integration pass to maximize contributor visibility before final synthesis.

The following imported block preserves expanded descriptive material from Person A's Onocoy case draft.

Onocoy is described as a decentralized Real-Time Kinematic (RTK) GNSS correction network. RTK correction is used to improve positioning precision from meter-level to centimeter-level ranges, with relevance for precision agriculture, robotics and automation, surveying, and industrial machinery control.
% TODO-CITE: [Person A sources: Teunissen and Montenbruck 2017; EUSPA 2023]

The draft emphasizes that Onocoy coordinates distributed reference stations operated by independent participants. These stations provide correction data that is aggregated and distributed through a protocol-mediated coordination layer.

The demand side is characterized as primarily professional and enterprise-oriented. In this view, demand for high-accuracy correction services is linked to operational workflows and continuity requirements rather than to purely speculative token narratives.
% TODO-CITE: [Person A sources: EUSPA 2023; Onocoy docs]

Provider requirements are also preserved explicitly. Reference-station operators are described as responsible for:
\begin{itemize}
    \item deploying compatible GNSS hardware,
    \item maintaining continuous uptime,
    \item maintaining calibration quality,
    \item providing stable connectivity.
\end{itemize}

Failure to maintain uptime or quality can reduce delivered service reliability and affect reward eligibility.
% TODO-CITE: [Person A sources: RTCM 2023; Onocoy docs]

Person A's token architecture framing is preserved as a dual-layer model:
\begin{itemize}
    \item ONO token for utility/governance and incentive-layer functions.
    \item Data Credits as non-transferable usage credits consumed during service access.
\end{itemize}

Users purchase Data Credits for service usage, and credits are consumed on use. This architecture is presented as an attempt to separate user-side payment flows from direct exposure to secondary-market token volatility.
% TODO-CITE: [Person A sources: Onocoy docs; Helium docs comparison]

The draft also notes that Onocoy remains exposed to broader crypto-market cycles and macro conditions despite utility-demand anchoring. Token volatility can still influence provider-side incentives and retention behavior.
% TODO-CITE: [Person A sources: Cong et al. 2021]

A further preserved analytical point is that Onocoy's central stress-relevant interaction is between relatively stable industrial demand and potentially unstable token-market dynamics. This interaction is treated as a key bridge to later stress-model parameterization and scenario testing.

\paragraph{Traceability note.} Full verbatim Person A submission snapshots are preserved in Appendix~\ref{subsec:contributor_draft_archive} for audit and attribution during drafting.
